\scanpage{64}
\begin{center}6 - 35 -\hfill n\textsuperscript{o} 392\end{center}

9) Soient $k$ un corps, $A$ l'anneau de polynômes $k[X,Y]$. Dans le $A$-module $M=A^2$, on considère la forme linéaire $f$ telle que $f(e_1)=X$, $f(e_2)=Y$, ($(e_1,e_2)$ étant la base canonique). Montrer que le noyau $L$ de $f$ est un $A$-module libre monogène, mais que le quotient $M/L$, (qui est isomorphe à un idéal de $A$), n'est pas réflexif.

10) Soit $A$ un anneau commutatif. Pour tout sous-module $R$ d'un $A$-module libre de type fini $L=A^n$, on désigne par $\underline c_1(R)$ l'idéal engendré par les $\langle x,x^*\rangle$, où $x$ parcourt $R$ et $x^*$ parcourt le dual $L^*$ ; on pose $\underline c_k(R)=\underline c_1(\operatorname{Im}(\bigwedge^k R))$, $\operatorname{Im}(\bigwedge^k R)$ étant l'image canonique de la puissance extérieure $k$-ème de $R$ dans $\bigwedge^k L$. Si $R_1\subset R_2$ sont deux sous-modules de $L$, on a $\underline c_k(R_1)\subset\underline c_k(R_2)$ pour tout $k$.

a) Soit $M$ un $A$-module de type fini : $M$ est isomorphe à un module quotient $L/R$, où $L=A^n$ pour un $n$ convenable. Montrer que la suite d'idéaux $(\underline a_k)_{k\geq0}$ telle que $\underline a_k^{\unsure{\prime}}=\underline c_{n-k}(R)$ pour $k\leq n$ et $\underline a_k=A$ pour $k>n$, est indépendante de l'expression de $M$ sous la forme $L/R$. (Considérer d'abord le cas où $L/R$ et $L'/R'$ sont isomorphes ; remarquer ensuite que $A^n/R$ est isomorphe à $A^{n+h}/(R\times A^h)$ pour tout $h>0$). On pose $\underline d_k(M)=\underline a_k$ pour tout $k\geq0$, et on dit que ce sont les \uline{idéaux déterminantiels} associés à $M$. On a $\underline d_k(M)\subset\underline d_{k+1}(M)$ pour $k\geq0$.

b) Si $A$ est un anneau principal et $M=L/R$ où $L$ est libre de type fini, montrer que les idéaux $\underline c_{k+1}(R)(\underline c_k(R))^{-1}$ sont les \uline{facteurs invariants} de $R$ dans $L$.

c) Soit $r=\unsure{\gamma|g}(M)$ le plus petit des cardinaux des systèmes de générateurs de $M$, et soit $r_0$ le plus petit entier $h$ tel que $\underline d_h(M)=A$. Montrer que $r_0\leq r$ et donner un exemple où $r_0<r$, (prendre pour $A$ un anneau de Dedekind).

d) Si $\underline a$ est l'annulateur de $M$, montrer que l'on a $\underline d_0(M)\subset\underline a$ et $\underline a^{r-k}\subset\underline d_k(M)$ pour $k\leq r=\unsure{\gamma|g}(M)$.

e) On suppose que $M$ est somme directe de sous-modules $M_i$, $(1\leq i\leq h)$. Montrer que l'on a :
\[
\underline d_k(M)=\sum\underline d_{k_1}(M_1)\ldots\underline d_{k_h}(M_h)
\]
où la somme est étendue aux suites finies $(k_i)_{1\leq i\leq h}$ telles que $\sum_{i=1}^h k_i=k$.

f) Si $N$ est un sous-module de type fini de $M$, on a $\underline d_k(M/N)\subset\underline d_k(M)$ pour tout $k\geq0$. Montrer que l'on a :
\[
\sum_{j=0}^k\underline d_j(N)\underline d_{k-j}(M/N)\subset\underline d_k(M).
\]
\margin{63}
